\chapter{Basic Parameter Estimation}

Parameter estimation is an important use case of conditional distribution. This section introduces only the basics of parameter estimation as a use case of conditional probability. More about parameter estimation can be found in control system relevant notebooks. 

\section{Posteriori Mean}

Let $\theta$ be the parameter to be estimated, and $x$ the measurement that reflects $\theta$ via measurement model $f(\theta, x)$ which is given in the form of joint PDF. Both $\theta$ and $x$ can be vectors. Under the scope of the discussion in this section, it is assumed that $f(\theta,x)$ is known. Without measurement $x$, the estimate of $\theta$ is given by
\begin{eqnarray}
	\hat{\theta} &=& \int_{-\infty}^{\infty} \theta f_\theta(\theta) d\theta \nonumber
\end{eqnarray}
where $f_\theta(\theta)$ is obtained from $f(\theta,x)$ using \eqref{eq:jointpdfdowngrade}. This is known as the \mync{priori estimation} of $\theta$. 

Given measurement $x$, the \mync{posteriori estimation} of $\theta$ can be obtained as follows. From \eqref{eq:conditionalpdf1a},
\begin{eqnarray}
	f_{\theta|X}(\theta|x) &=& \dfrac{f_{X|\theta}(x|\theta)f_\theta(\theta)}{f_X(x)} \label{eq:conditionalpdf2} \\
	\hat{\theta} &=& \int_{-\infty}^{\infty} \theta f_{\theta|X}(\theta|x) d\theta \label{eq:posterioriestimation}
\end{eqnarray}
where $f_{X|\theta}(x|\theta)$ is known as the \mync{likelihood function} that describes the likelihood of measuring $x$ if the actual parameter(s) is $\theta$. Distribution $f_X(x)$ is known as the \mync{evidence}. With fixed $x$, $f_X(x)$ becomes a constant.

Equation \eqref{eq:conditionalpdf2} is essentially
\begin{eqnarray}
	\textup{posteriori} &=& \dfrac{\textup{likelihood}\times\textup{priori}}{\textup{evidence}} \label{eq:posteriorimemo}
\end{eqnarray} 

Notice that all the components in \eqref{eq:conditionalpdf2} and \eqref{eq:posteriorimemo} are derived from the joint distribution $f(\theta, x)$. With known analytical $f(\theta, x)$ and a set of measurement $x$, \eqref{eq:conditionalpdf2} can be used to calculate the estimate. 

The use of posteriori estimation \eqref{eq:posteriorimemo} assumes that $f(\theta, x)$ is know, from where the likelihood, the priori and the evidence can be derived given a specific measurement $x$. However, this is often not true in reality, where $f(\theta, x)$ is often unknown and the likelihood, the priori and the evidence are either unknown or only partially known. In the next Section \ref{sec:mlemap}, \myabb{Maximum Likelihood Estimation}{MLE} is introduced as an example of alternation to the posteriori estimation that does not require the full picture of $f(\theta, x)$.

\section{MLE and MAP} \label{sec:mlemap}

Alternative to \eqref{eq:posterioriestimation}, there are at least two other types of commonly seen estimations, namely \mync{Maximum Likelihood Estimation}[MLE] and \mync{Maximum A Posteriori}[MAP] estimations. Details of them can be found in control system relevant notebooks. A brief introduction is given below.

\subsection{MLE}

In many occasions, the joint distribution $f(\theta, x)$ is unknown. Instead, the measurement $x$ can be formulated as a function of $\theta$ plus noise. This function is known as the measurement model or output model of the system. In this case, the likelihood $f_{X|\theta}(x|\theta)$ can be derived from the measurement model, whereas the priori estimation and evidence are unknown. As a result, \eqref{eq:conditionalpdf2} and \eqref{eq:posterioriestimation} cannot be used to calculate the estimate.

In this case, consider maximizing the likelihood function to obtain the estimate as follows.
\begin{eqnarray}
	\hat{\theta} &=& \argmax_\theta f_{X|\theta}(x|\theta) \nonumber
\end{eqnarray}
This is known as MLE. MLE is useful when only the measurement model is known.

There are many ways to solve an MLE problem. In the special cases where $f_{X|\theta}(x|\theta)$ is a special distribution, such as Gaussian and Laplace, the MLE becomes \mync{Weighted Least Squares}[WLS] estimator and \mync{Least Absolute Value}[LAV] estimator, respectively.

A general way of solving MLE is given below. Define the cost function as
\begin{eqnarray}
	J &=& -\sum_{i=1}^{m}\textup{ln}f_{X|\theta}(x|\theta) \nonumber 	
\end{eqnarray}
and use 
\begin{eqnarray}
	\hat{\theta} &=& \arg\min_{x} J \nonumber
\end{eqnarray}
which can be solved using commonly used optimization algorithms such as Newton-Raphson method.

\subsection{MAP}

MAP estimation is also known as the Bayesian estimation. Like the posteriori mean method, MAP also calculates the posteriori distribution \eqref{eq:conditionalpdf2}. But instead of \eqref{eq:posterioriestimation}, the estimate is calculated using
\begin{eqnarray}
	\hat{\theta} &=& \argmax_\theta f_{\theta|X}(\theta|x) \nonumber \\
	&=&  \argmax_\theta \dfrac{f_{X|\theta}(x|\theta)f_\theta(\theta)}{f_X(x)} \label{eq:mapestimate1}
\end{eqnarray}

Notice that \eqref{eq:mapestimate1} is equivalent to
\begin{eqnarray}
	\hat{\theta} &=& \argmax_\theta f_{X|\theta}(x|\theta)f_\theta(\theta) \label{eq:mapestimate2}
\end{eqnarray}
as $f_X(x)$ is a constant for any $x$ and does not affect the result of \eqref{eq:mapestimate1}. In this sense, comparing with the posteriori mean method, MAP does not require $f_X(x)$ nor the joint distribution $f(\theta, x)$ in the calculation. Comparing with MLE, MAP further requires the priori $f_\theta(\theta)$.

Equation \eqref{eq:mapestimate2} can be solved as follows.
\begin{eqnarray}
	J &=& -\sum_{i=1}^{m}\textup{ln}f_{X|\theta}(x|\theta) - \textup{ln}f_\theta(\theta) \nonumber \\
	\hat{\theta} &=& \arg\min_{x} J \nonumber
\end{eqnarray}
which can likewise be solved using commonly used optimization algorithms.

\begin{mdframed}
	\noindent \textbf{MLE versus MAP versus Posteriori Mean}
	
	It is clear that MLE require the least information among the three to carry out the calculation. MAP requires more information. Posteriori mean method requires the most information.
	
	MLE maximizes the likelihood function $f_{X|\theta}(x|\theta)$ only. MAP maximizes the product of likelihood function and the priori estimation, $f_{X|\theta}(x|\theta)f_\theta(\theta)$. From that sense, MLE can be taken as a special case of MAP where $f_\theta(\theta)$ is constant for all $\theta$.
	
\end{mdframed}

\section{Parameter Estimation Benchmark}

Estimation error is almost always unavoidable due to the limited size of the samples. However, it is desirable that with an appropriate estimation method, the estimates are unbiased from the true value of the population, i.e., the expectation of the estimates should equal the corresponding values of the population. In that case, the estimation is known as \mync{unbiased estimation}. 

Consider unbiased estimation. The accuracy of an estimation method can be evaluated as follows.

The maximum information contained in the samples can be described by \mync{Fisher Information Matrix} named after Sir Ronald Fisher. The following equation describes the relation between $\textup{Var}(\hat{\theta})$ and Fisher Information Matrix $I(\theta)$ for a scalar unbiased estimator.
\begin{eqnarray}
	\textup{Var}(\hat{\theta}) &\geq& I(\theta)^{-1} \label{eq:crlb}
\end{eqnarray}
which is known as the \mync{Cram\'{e}r–Rao Lower Bound}[CRLB] of the estimator. From \eqref{eq:crlb}, it can be seen that the ``larger'' $I(\theta)$, the lower $\textup{Var}(\hat{\theta})$  can possibly reach.

It is possible (although in many cases difficult) to check and prove whether an estimator hits the bound and provides the smallest possible $\textup{Var}(\hat{\theta})$ by comparing it with its corresponding CRLB. If its variance equals to CRLB, the estimator an \mync{optimal estimator}. A more general CRLB for biased estimator can also be derived. The detailed derivations of Fisher Information Matrix and general CRLB are not given in this notebook.

The magnitude of the estimation error can be evaluated using \mync{Mean Squared Error}[MSE] and \mync{Root Mean Squared Error}[RMSE] as defined below. 

Let $\theta$ be a parameter and $\hat{\theta}_i$ its estimation in the $i$th Monte-Carlo run. Let $n$ be the total number of Monte-Carlo runs. Define MSE and RMSE over the $n$ Monte-Carlo runs as follows.
\begin{eqnarray}
	\textup{MSE} &=& \dfrac{1}{n}\sum_{i=1}^{n}\left(\hat{\theta}_i-\theta\right)^2 \nonumber \\
	\textup{RMSE} &=& \sqrt{\textup{MSE}} \nonumber
\end{eqnarray}
Notice that in practice, MSE and RMSE are meaningful only when $n$ is large enough so that their values converge. Let $n$ approaches infinity and we can observe that
\begin{eqnarray}
	\textup{Var}(\hat{\theta}) = E\left[\left(\hat{\theta}-E\left[\hat{\theta}\right]\right)^2\right] = E\left[\left(\hat{\theta}-\theta\right)^2\right] = \textup{MSE} \nonumber
\end{eqnarray}
for unbiased estimation because $E\left[\hat{\theta}\right]=\theta$, and the MSE approaches the variance of the estimate. From this point of view, we can use the variance of the estimate, $\textup{Var}(\hat{\theta})$, interchangeably with MSE as an evaluation of the performance of the estimation. In most literatures, $\textup{Var}(\hat{\theta})$ is used, implying that the result is theoretical and derived under the assumption of ``infinite runs''.

The smaller $\textup{Var}(\hat{\theta})$, the more precise the estimation. There is a limitation to the minimum value of $\textup{Var}(\hat{\theta})$ due to CRLB. This is intuitive because there is only so much information contained in the samples, given limited sample size and measurement error. So long as the estimation is done using those samples, its precision is limited. 

